% !TeX program = xelatex
% ============================================================
% Аннотация — документ КТ1 исследовательской (академической) ВКР
% ОП «Программная инженерия» ФКН НИУ ВШЭ.
%
% ЖАНР (Инструкция по подготовке документов к защите тем ВКР 2025/26 +
% ГОСТ 7.32-2017 §5.3): набор ПОЛУЖИРНЫХ ЛЕЙБЛОВ — структурных частей
% аннотации — со значением следом. НЕ ЕСПД: без штампа, без кода
% документа, без листа утверждения, без содержания и без нумерованных
% разделов. Требования КТ1: аннотация СО СПИСКОМ ИСТОЧНИКОВ — не менее
% 7–8 страниц; источников — не менее 5; ТИТУЛЬНЫЙ ЛИСТ ПОДПИСЫВАЮТ
% студент И руководитель (подписи размещаются в разделе «Подписи»).
%
% Как пользоваться: замените жёлтые \hseFill{…} своим текстом. Если по
% какой-то структурной части сведений нет — часть опускается, при этом
% ПОСЛЕДОВАТЕЛЬНОСТЬ остальных сохраняется (инструкция §3). Данные о
% себе, руководителе и годе задаются в настройках проекта (макросы \hse…).
% ============================================================
\documentclass[12pt,a4paper]{extarticle}
% !TeX program = xelatex
% ============================================================
% common/annotation_preamble.tex — преамбула Аннотации (документ КТ1
% исследовательской ВКР ОП «Программная инженерия»).
%
% ЭТО НЕ ЕСПД-документ: без штампа/рамки, без кода документа, без листа
% утверждения и без содержания. Оформление по ГОСТ 7.32-2017: поля
% 30/15/20/20, Times New Roman 12 пт, интервал 1,5, обычный номер
% страницы снизу. Данные — из common/meta (\hse-макросы). Ссылки [n]
% печатаются В СТРОКУ, список источников — по ГОСТ Р 7.0.5-2008.
% ============================================================

% ── Шрифты (автогенерируются из настроек проекта) ───────────
\newcommand{\hseDefaultMono}{Courier New}
\input{../common/fonts}
\usepackage[english,russian]{babel}

% ── Геометрия по ГОСТ 7.32-2017 §6.1.1 (без места под штамп) ─
\usepackage[a4paper,left=30mm,right=15mm,top=20mm,bottom=20mm]{geometry}
\usepackage{setspace}
\onehalfspacing
\usepackage{indentfirst}
\setlength{\parindent}{1.25cm}
\setlength{\parskip}{6pt}

% ── Типографика и таблицы ───────────────────────────────────
\usepackage{microtype}
\usepackage{csquotes}
\usepackage{enumitem}
\setlist[itemize,1]{label=--}
\usepackage{graphicx}
\usepackage[table]{xcolor}
\usepackage{array}
\usepackage{tabularx}
\usepackage{booktabs}
\usepackage{float}

% ── URL и переносы ──────────────────────────────────────────
\usepackage{url}
\usepackage{xurl}
\urlstyle{same}
\usepackage{hyphenat}

% ── Цитирование: ссылки [n] В СТРОКУ (как в реальных аннотациях),
%    список источников — ГОСТ Р 7.0.5-2008. НЕ переопределяем \@cite
%    в надстрочный вид (в отличие от ЕСПД-преамбулы). ────────────
\usepackage{cite}
\makeatletter
\renewcommand\@biblabel[1]{#1.}
\makeatother

\usepackage{lastpage}
\usepackage{hyperref}
\hypersetup{
    unicode=true,
    breaklinks=true,
    colorlinks=true,
    linkcolor=black,
    urlcolor=blue,
    citecolor=blue,
    pdfborder={0 0 0}
}

% ── Данные проекта (\hse…), затем «человеческие» имена (commands) ──
% meta даёт \hseFill/\hseOptional и все \hse-макросы; commands —
% \signdateblank, \hseExecutorsBlock-зависимости (\authorsignatureline,
% \studentgroup, \studentname) для блока «Выполнил» в титуле.
\input{../common/meta}
\input{../common/commands}

% Без штампа: обычный номер страницы снизу по центру (ГОСТ 7.32 §6.3.1).
\pagestyle{plain}

\begin{document}
\sloppy
% ============================================================
% common/annotation_title.tex — титульный лист Аннотации ВКР ОП ПИ
% по официальной форме «Титул Аннотация защита темы ВКР ПИ».
%
% Лёгкий титул БЕЗ штампа/рамки/листа утверждения: УДК, две колонки
% СОГЛАСОВАНО (руководитель проекта) + УТВЕРЖДАЮ (академический
% руководитель ОП), затем «Аннотация выпускной квалификационной
% работы», тема, направление и блок «Выполнил» (в командной работе —
% состав исполнителей через \hseExecutorsBlock).
% ============================================================
\begin{titlepage}
\thispagestyle{empty}
\singlespacing
\footnotesize
\setlength{\parindent}{0pt}
\setlength{\parskip}{0pt}

\begin{center}
\textbf{ПРАВИТЕЛЬСТВО РОССИЙСКОЙ ФЕДЕРАЦИИ}\\
\textbf{ФЕДЕРАЛЬНОЕ ГОСУДАРСТВЕННОЕ АВТОНОМНОЕ}\\
\textbf{ОБРАЗОВАТЕЛЬНОЕ УЧРЕЖДЕНИЕ ВЫСШЕГО ОБРАЗОВАНИЯ}\\
\textbf{НАЦИОНАЛЬНЫЙ ИССЛЕДОВАТЕЛЬСКИЙ УНИВЕРСИТЕТ}\\
\textbf{«ВЫСШАЯ ШКОЛА ЭКОНОМИКИ»}

\vspace{0.35cm}
\hseFaculty\\
Образовательная программа «\hseSpecialization»
\end{center}

\vspace{0.25cm}
\noindent
УДК~\hseUDC

\vspace{0.4cm}
\noindent
\begin{minipage}[t]{0.46\textwidth}
\raggedright
\textbf{СОГЛАСОВАНО}\\
Руководитель проекта

\vspace{0.15cm}
\strut\hseSupervisorTitle\\[-0.05cm]
\rule{6cm}{0.4pt}\\[-0.05cm]
{\scriptsize Должность, место работы, ученая степень}

\vspace{0.12cm}
\strut\hseSupervisorName\\[-0.05cm]
\rule{6cm}{0.4pt}\\[-0.05cm]
{\scriptsize Подпись, И.О. Фамилия}

\vspace{0.1cm}
\signdateblank
\end{minipage}
\hfill
\begin{minipage}[t]{0.46\textwidth}
\raggedright
\textbf{УТВЕРЖДАЮ}\\
\strut\hseAcademicSupervisorTitle

\vspace{0.35cm}
\mbox{\rule{6cm}{0.4pt}}\\[-0.05cm]
\hseAcademicSupervisorName

\vspace{0.12cm}
«\rule{0.8cm}{0.4pt}» \rule{2.5cm}{0.4pt} \hseYear~г.
\end{minipage}

\vspace{0.95cm}
\begin{center}
\textbf{\normalsize Аннотация}

\vspace{0.12cm}
\textbf{\normalsize выпускной квалификационной работы}

\vspace{0.35cm}
на тему
\underline{\makebox[0.72\textwidth][c]{\hseProjectName}}

\vspace{0.3cm}
по направлению подготовки бакалавров \hseProgramCode\ «\hseSpecialization»
\end{center}

\vspace{0.95cm}
\noindent\hfill
\begin{minipage}[t]{0.52\textwidth}
\raggedright
% Командная работа (shared-база) — состав исполнителей; solo/личная — один.
\ifhseSharedDoc
\hseExecutorsBlock
\else
Выполнил\\
студент группы \hseGroup\\
образовательной программы\\
\hseProgramCode\ «\hseSpecialization»

\vspace{0.25cm}
\strut\hseAuthorName\\[-0.05cm]
\rule{6cm}{0.4pt}\\[-0.05cm]
{\scriptsize И.О. Фамилия}

\vspace{0.15cm}
\rule{3.6cm}{0.4pt}\quad\rule{2.2cm}{0.4pt}\\[-0.05cm]
{\scriptsize Подпись, Дата}
\fi
\end{minipage}

\vfill
\begin{center}
\textbf{\hseCity{} \hseYear}
\end{center}
\end{titlepage}
\clearpage



% ============================================================
%  ENGLISH ANNOTATION BODY (project.lang == en)
%  §2.5.3.1 / §2.3.4.2: an English research VKR is written in English.
%  A Russian AND an English abstract are BOTH required, but this light
%  document is the annotation itself — here it is rendered in English.
%  References follow IEEE style. Fill the yellow \hseFill{…} placeholders.
% ============================================================
\begin{center}
  \textbf{\large Annotation}
\end{center}

\vspace{0.4em}

% ── Keywords (GOST 7.32-2017 §5.3.2.1 / §6.12.2): 5–15 words or
%    phrases, in UPPERCASE, in the nominative case, on one line separated
%    by commas, without a full stop at the end. ──
\noindent\textbf{Keywords:} \hseFill{5–15 WORDS IN UPPERCASE, NOMINATIVE CASE, SEPARATED BY COMMAS, NO FULL STOP AT THE END}

\medskip
\noindent\textbf{Assessment of the current state of the problem being solved.} The problem under consideration is~\hseFill{state the scientific problem briefly}. Today it is addressed with \hseFill{list existing approaches, methods, works} \cite{example-en}. Their main limitations are: \hseFill{limitation 1}; \hseFill{limitation 2}~--- which is what defines the direction of the present research.

\medskip
\noindent\textbf{Basis and initial data for the research.} The work is based on the curriculum of the «\hseSpecialization» programme and the thesis topic approved by the HSE FCS order. Initial data: \hseFill{existing groundwork, data, publications and software the work builds upon}.

\medskip
\noindent\textbf{Relevance and novelty of the research topic.} The relevance stems from \hseFill{the practical and scientific significance of the task}. The scientific novelty of the work lies in~\hseFill{what exactly is proposed for the first time}.

\medskip
\noindent\textbf{Object of study:} \hseFill{what is studied as a whole}.

\smallskip
\noindent\textbf{Subject of study:} \hseFill{the specific aspect of the object on which the work focuses}.

\smallskip
\noindent\textbf{Aim of the research:} \hseFill{a single formulation of the scientific result to be achieved}.

\medskip
\noindent\textbf{Research objectives:}
\begin{enumerate}[label=\arabic*), leftmargin=1.3cm, topsep=2pt, itemsep=2pt]
  \item \hseFill{carry out a review of existing approaches};
  \item \hseFill{propose and formalise a method (model, algorithm)};
  \item \hseFill{implement a software tool for the experiments};
  \item \hseFill{run a computational experiment and evaluate the results};
  \item \hseFill{draw conclusions on the applicability}.
\end{enumerate}

\medskip
\noindent\textbf{Research methods:} \hseFill{list the methods: analytical review, formal modelling, computational experiment, statistical analysis, etc.}

\medskip
\noindent\textbf{Validity of the scientific results} is confirmed by \hseFill{reproducibility of the experiments, comparison with existing approaches, statistical significance of the differences}.

\medskip
\noindent\textbf{Expected results and their novelty:}
\begin{enumerate}[label=\arabic*), leftmargin=1.3cm, topsep=2pt, itemsep=2pt]
  \item \hseFill{method / model / algorithm};
  \item \hseFill{software implementation for running the experiments};
  \item \hseFill{experimental evaluation and comparison with baseline methods};
  \item \hseFill{conclusions on the limits of applicability}.
\end{enumerate}

\medskip
\noindent\textbf{Field of application of the results:} \hseFill{specify the field of application}.

\medskip
\noindent\textbf{Expected adoption of the results.} \hseFill{where and in what form the results are expected to be adopted: a product, a service, a publication, open-source code}.

\medskip
\noindent\textbf{Economic efficiency or significance of the work.} The results of the work can be used \hseFill{where and by whom the obtained results are applicable; what the practical or economic gain is}.

\medskip
\noindent\textbf{Forecast assumptions on the development of the object of study.} \hseFill{how the object of study may evolve and how the work may be continued}.

\bigskip
% ── References in IEEE style (§2.5.3.1): for an English research VKR the
%    reference list and in-text citations follow IEEE style — numbered
%    [n], author initials first, article titles in quotes,
%    venue/volume/number/pages, accessed date for online sources. Every
%    source must be cited with \cite{…} in the text. ──
\noindent\textbf{References}
\vspace{-0.3em}
\renewcommand{\refname}{}
\begin{thebibliography}{99}
  \bibitem{example-en}
  A.~B.~Author and C.~D.~Coauthor, \enquote{Title of the paper,} in \emph{Proc. of the Conference}, New York, NY, USA: ACM, 2024, pp.~10--20.

  \bibitem{example-web}
  Resource Name. (2026). Accessed: Jan.~1, 2026. [Online]. Available: \url{https://example.org}
\end{thebibliography}



\end{document}
